%=============================
% ENTORNOS PERSONALIZADOS

% Cuadro ejemplo
\newcounter{example}[section]
\renewcommand{\theexample}{\thesection.\arabic{example}}

\newtcolorbox{ejemplo}[1][]{
    enhanced,
    breakable,
    colback=gray!12,
    boxrule=0pt,
    frame hidden,
    borderline west={2pt}{0pt}{black},
    sharp corners,
    left=8pt,
    right=8pt,
    top=6pt,
    bottom=6pt,
    fonttitle=\bfseries,
    title={Ejemplo~\theexample},
    coltitle=black,
    code={\refstepcounter{example}},
    #1
}

% Notas
\newtcolorbox{nota}{
  enhanced,
  breakable,
  colback=gray!6,
  boxrule=0pt,
  frame hidden,
  borderline west={3pt}{0pt}{gray!60},
  sharp corners,
  left=8pt,
  right=8pt,
  top=6pt,
  bottom=6pt,
  fontupper=\small
}

% Ejercicios 
\newtcolorbox{ejercicio}[1][]{
  enhanced,
  breakable,
  colback=white,
  colframe=black!70,      % Color del marco y del título
  colbacktitle=black!70, % Fondo del título
  coltitle=white,        % Texto del título en blanco
  fonttitle=\bfseries\small,
  title={EJERCICIO #1},   % El título es "EJERCICIO" + lo que tú escribas
  detach title,
  % Diseño del encabezado
  before upper={\textsf{\textbf{\textcolor{black!70}{EJERCICIO #1}}}\smallskip\hrule\medskip\par},
  boxrule=0pt,
  frame hidden,
  borderline north={1.5pt}{0pt}{black!70}, % Línea superior gruesa
  borderline south={0.5pt}{0pt}{black!20}, % Línea inferior sutil
  sharp corners,
  left=0pt,   % Alineado a la izquierda sin sangría
  right=0pt,
  top=5pt,
  bottom=5pt,
  fontupper=\small
}

%=============================